\section[Summary]{Summary}

\begin{frame}{Integrated Textbook Practice}
\small
\begin{columns}[T]
\column{.48\textwidth}
\textbf{Invertibility}
\[
A=\begin{bmatrix}2&3&4\\2&3&4\\2&3&4\end{bmatrix}.
\]
What is the cheapest reason that $A$ is singular?

\medskip
\textbf{Block multiplication}
If $X=[\,X_1\ X_2\,]$, write $X^TX$ in block form.

\column{.48\textwidth}
\textbf{LU practice}
\[
A=\begin{bmatrix}
2&-4&-2&3\\
6&-9&-5&8\\
2&-7&-3&9\\
4&-2&-2&-1\\
-6&3&3&4
\end{bmatrix}.
\]
The matrix has three pivot columns. Which parts of $L$ come from elimination, and which parts come from $I_5$?
\end{columns}
\end{frame}

\begin{frame}{Practice Connections}
\small
\begin{columns}[T]
\column{.48\textwidth}
The first matrix has identical rows, so it has fewer than three pivot positions. By the Invertible Matrix Theorem, it is singular.

\medskip
For $X=[\,X_1\ X_2\,]$,
\[
X^TX=
\begin{bmatrix}
X_1^TX_1&X_1^TX_2\\
X_2^TX_1&X_2^TX_2
\end{bmatrix}.
\]

\column{.48\textwidth}
For the LU practice problem, elimination supplies the first three columns below the diagonal of $L$. The remaining columns are chosen from $I_5$ so that $L$ is unit lower triangular.

\medskip
The same row operations that reduce $A$ to $U$ must reduce $L$ to $I_5$.
\end{columns}
\end{frame}

\begin{frame}{Summary}
\begin{enumerate}
  \item In a matrix equation, multiplication order determines which inverse goes on which side.
  \item For a square matrix, pivots, solution behavior, column properties, and transformation properties describe the same invertibility condition.
  \item A valid block product requires matching inner block dimensions.
  \item An LU factorization reuses elimination work when the coefficient matrix stays fixed.
\end{enumerate}
\medskip
\begin{alertblock}{Exit check}
For each topic, identify the calculation that reveals the structure: row reduction, block multiplication, or triangular solution.
\end{alertblock}
\end{frame}
